\section{Two finite initializations}\label{sec:initialization}

The doubling argument requires a spatial purity at length $n$ and a physical
purity at length $2n$, at the same inverse temperature.  We establish two such
pairs.  The first begins at length $60$; the second begins at length $2304$ and
gives the stronger eventual gap.  In both cases the inputs are a few short-chain
partition functions and one periodic trial vector.  This section explains how
those finite inputs imply the required purities.

\subsection{Thermal data and elementary moment bounds}

We use the sectors from Lemma~\ref{lem:sectors}: $J_n=I_n+2O_n$ is the
$D_2$-invariant moment, $I_n$ is the moment fixed also by the cyclic rotation,
and $N_n$ is the moment in one nontrivial $D_2$ sector.  Each $O_n$ occurs twice
and each $N_n$ occurs three times in the full trace.  In particular,
\begin{equation}\label{eq:sector-inversion}
\begin{aligned}
 J_n&=\frac{Z_n+3Z_n(P)}4,&
 N_n&=\frac{Z_n-Z_n(P)}4,\\
 I_n&=\frac{Z_n+3Z_n(P)+8Z_n(C)}{12},&
 O_n&=\frac{Z_n+3Z_n(P)-4Z_n(C)}{12}.
\end{aligned}
\end{equation}
The inverse temperature is common to every term of each identity.

\begin{lemma}[Finite thermal inputs]\label{lem:thermal-inputs}
For $a=700741/500000$, the following tables give centers for the shifted
partition functions $Z_n(b,g)$ in units $10^{-6}$.
The absolute errors are less than $10^{-5}$ at $b=21/2$ and less than
$30\cdot10^{-6}$ at $b=49/4$.
\[
\begin{array}{c|rrr}
 &\multicolumn{3}{c}{b=21/2}\\
 n&g=1&g=P&g=C\\ \hline
 6&8945303&346734&939987\\
 8&3741281&567167&984108\\
 10&2327547&721090&1002613
\end{array}
\qquad
\begin{array}{c|rr}
 &\multicolumn{2}{c}{b=49/4}\\
 n&g=1&g=P\\ \hline
 4&124885379&65757\\
 5&2925&2960742\\
 6&12870614&286599\\
 7&54069&1920414\\
 8&4639283&510140\\
 9&195455&1514846\\
 10&2654823&675419\\
 11&377629&1314544\\
 12&1900547&787405
\end{array}
\]
\end{lemma}

Appendices~\ref{sec:thermal120} and~\ref{sec:thermal72} prove the respective
tables by rational approximations to the physical matrix exponentials,
with explicit error bounds.  Lemma~\ref{lem:trial-inputs} supplies the other
finite input:
\begin{equation}\label{eq:trial-thermal-consequence}
 E_0(72)<-72a,\qquad E_0(120)<-120a.
\end{equation}
Consequently $Z_{72}(b)>1$ and $Z_{120}(b)>1$ for every $b>0$: in each case
one Boltzmann factor already exceeds one.

Here are two elementary ways to extract spectral information from the tables.

\begin{lemma}[Capped masses]\label{lem:cappedmass}
Let $(t_i)$ be a finite or countable family with $0\le t_i\le h$ and
$\sum_i t_i\le m$, where $h>0$ and $m\ge0$.  For $r\ge1$, put
\begin{equation}\label{eq:capped-function}
 \mathcal C(m,h,r)=q h^r+(m-qh)^r,\qquad q=\lfloor m/h\rfloor.
\end{equation}
Then $\sum_i t_i^r\le\mathcal C(m,h,r)$.  Without the individual cap,
$\sum_i t_i^r\le m^r$.
\end{lemma}

\begin{proof}
For a finite family at fixed total mass, transfer mass from the smaller of two
coordinates in $(0,h)$ to the larger until one reaches $0$ or $h$.  Convexity
shows that the sum of $r$th powers does not decrease.  Repetition leaves at most
one coordinate in $(0,h)$, giving~\eqref{eq:capped-function} at the actual
total mass.  That expression is continuous and nondecreasing in the total
mass, so it is bounded by its value at $m$.  Apply this argument to finite
subfamilies and take their increasing union for a countable family.  The
uncapped assertion follows directly from
$t_i^r\le t_i(\sum_jt_j)^{r-1}$ when the total mass is positive; the zero case
is immediate.
\end{proof}

\begin{lemma}[Polynomial filter]\label{lem:polynomial-filter}
Let $\epsilon\ge0$, $v>0$, and let $(\lambda_i)$ be a real eigenvalue list with moments
$M_j=\sum_i\lambda_i^j$, and suppose $|M_j-c_j|\le\epsilon$ for the
indices occurring in a polynomial $F(x)=\sum_j f_jx^j$.
Assume these moments converge absolutely and $F\ge0$ on $\R$.
If
\[
 B=\sum_j f_jc_j+\epsilon\sum_j|f_j|
 \quad\hbox{and}\quad F(x)>B\ \hbox{for }|x|\ge v,
\]
then every $|\lambda_i|<v$.
\end{lemma}

\begin{proof}
Absolute convergence gives $\sum_iF(\lambda_i)\le B$.  Since every summand
is nonnegative, one eigenvalue with $|\lambda_i|\ge v$ would contradict this
bound.
\end{proof}

All moments used below converge absolutely by Proposition~\ref{prop:transfer}.
The coefficient sums in~\eqref{eq:sector-inversion} are at most one in
absolute value, so the thermal error bound also bounds each sector-moment
error.  This observation includes the odd moments used in the second
initialization; their signs are retained in the polynomial filter.

\subsection{The initialization at length sixty}

\begin{proposition}\label{prop:seed60}
At inverse temperature $b=105/4$,
\[
 S(60,b)>\frac78,\qquad T(120,b)>\frac78.
\]
\end{proposition}

\begin{proof}
We first bound sector moments and then increase the inverse temperature.
The trial energy forces a dominant spatial weight, giving the spatial
purity; its physical Boltzmann factor will give the physical purity
separately.

Begin at the lower inverse temperature $b_0=21/2$.  From the $b_0=21/2$ table in
Lemma~\ref{lem:thermal-inputs} and~\eqref{eq:sector-inversion}, upper bounds
for $I_{10},O_{10},N_{10}$ are, respectively,
\[
 m_I=\frac{1042655}{10^6},\qquad
 m_O=\frac{40041}{10^6},\qquad
 m_N=\frac{401625}{10^6}.
\]
For clarity, their margins over the sector centers plus $10^{-5}$ are
$19/(12\cdot10^6)$, $7/(12\cdot10^6)$ and $3/(4\cdot10^6)$.

Apply Lemma~\ref{lem:polynomial-filter} to
\[
 F_I(x)=x^6(-2+10x^2)^2,\qquad
 F_N(x)=x^6(-25+100x^2)^2.
\]
For a filter $x^6(c_0+c_1x^2)^2$, its summed upper bound is the centered
combination of moments $6,8,10$ with coefficients $c_0^2,2c_0c_1,c_1^2$,
plus $10^{-5}(|c_0|+|c_1|)^2$.  Direct rational evaluation gives
\[
 |\lambda|<U:=\frac{100245}{100000}\quad\hbox{in }I,
 \qquad
 |\lambda|<V:=\frac{88643}{100000}\quad\hbox{in }N.
\]
Indeed the filter value at the indicated endpoint exceeds its summed upper
bound by more than $0.0301$ for $I$ and $0.3463$ for $N$.
Each filter is strictly increasing with $|x|$ beyond its endpoint: there
$c_1>0$ and $c_0+c_1x^2>0$.

The rational inequalities
$U^{10}<m_I<2U^{10}$ and $V^{10}<m_N<2V^{10}$ allow the two-piece form
of Lemma~\ref{lem:cappedmass}.  For $k=30,120$ define
\[
 A_k=U^k+(m_I-U^{10})^{k/10},\qquad
 B_k=m_O^{k/10},\qquad
 C_k=V^k+(m_N-V^{10})^{k/10}.
\]
They bound $I_k,O_k,N_k$, respectively.  Because these are even moments,
all are nonnegative, and hence
\begin{equation}\label{eq:short-moment-upper}
 Z_k(b_0)\le A_k+2B_k+3C_k,\qquad
 Z_k(b_0,P),\ Z_k(b_0,C)\le A_k+2B_k.
\end{equation}
Put $m=A_{30}+2B_{30}+3C_{30}$, $p=A_{30}+2B_{30}$, and
$q=A_{120}+2B_{120}+3C_{120}$.  With
\[
 D=\frac{1450}{1000},\quad E=\frac{1202}{1000},\quad
 \tau=\frac{1225}{1000},\quad\tau_0=\frac{1264}{1000},\quad
 R_q=\frac{684}{10000},
\]
the remaining scalar estimates are
\begin{equation}\label{eq:seed60-scalars}
\begin{gathered}
 m^5<D^2,\qquad p^5<E^2,\qquad
 q\ge1,\qquad (q-1)^5<R_q^2,\\
 \frac{D+11E}{12}\le\tau,\qquad
 \frac{D+3E}{4}\le\tau_0.
\end{gathered}
\end{equation}
For example, the positive margins in the first, second and fourth strict
inequalities exceed $0.001068$, $0.000433$ and $0.0000472$, respectively.
Every number in these comparisons is rational and is specified above.

We now increase the inverse temperature from $b_0$ to
$b=(5/2)b_0$.  For positive Boltzmann factors $z_i$,
$\sum_i z_i^{5/2}\le(\sum_i z_i)^{5/2}$.  Thus
\eqref{eq:short-moment-upper}--\eqref{eq:seed60-scalars} imply
\[
 Z_{30}(b)\le D,\qquad Z_{30}(b,P),Z_{30}(b,C)\le E.
\]
By~\eqref{eq:sector-inversion}, the $I$ and $J$ portions of the thirtieth
spatial moment are at most $\tau$ and $\tau_0$.

Let $w_i=|\lambda_i(X_b)|^{30}$.  Then $\sum_iw_i\le D$, whereas
\eqref{eq:trial-thermal-consequence} gives $\sum_iw_i^4=Z_{120}(b)>1$.
Since $D(0.88)^3<1$, some weight $w_0>0.88$ exists.  It is unique because
$2(0.88)>D$.  It belongs to $I$: a weight in $O$ would occur at least twice,
and one in $N$ at least three times.  These statements count all
multiplicities, including equal moduli of opposite-sign spatial eigenvalues.

To sharpen the lower bound on $w_0$, suppose $0\le t\le0.998$ and $w_0\ge t$.
Denote the remaining
$I$ mass by $x$, the full $O$-pair mass by $y$, and the full $N$-triple mass
by $z$.  Their squared weights sum to at most $x^2+y^2/2+z^2/3$, and
\[
 0\le x\le\tau-t,\qquad 0\le y,\qquad
 x+y\le\tau_0-t,\qquad z\le D-t-x-y.
\]
The divisors $2$ and $3$ use the identical weight lists in the repeated
sectors.  Since $D>\tau_0$, the last upper bound is nonnegative on the entire
displayed polygon.  The convex quadratic
$x^2+y^2/2+(D-t-x-y)^2/3$ has its maximum at one of its four vertices.  Write
\begin{equation}\label{eq:residual-R}
 R(t)=\max_{(x,y)\in\mathcal V_t}
 \left(x^2+\frac{y^2}{2}+\frac{(D-t-x-y)^2}{3}\right),
\end{equation}
where
\[
 \mathcal V_t=\{(0,0),(\tau-t,0),
 (\tau-t,\tau_0-\tau),(0,\tau_0-t)\}.
\]
Thus the remaining squared weights sum to at most $R(t)$, and their fourth
powers sum to at most $R(t)^2$.  Exact evaluation gives
\[
\begin{array}{c|ccc}
 t&0.88&0.99&0.998\\ \hline
 R(t)&0.1359&0.0721&0.068404
\end{array}
\]
and
\[
 1-0.99^4-R(0.88)^2=0.02093518,\qquad
 1-0.998^4-R(0.99)^2=0.002777621984.
\]
If $w_0\le0.99$, the first equality would contradict
$\sum_iw_i^4>1$; the second then excludes $w_0\le0.998$.
Keeping the dominant contribution in the fourth moment yields
\[
 S(60,b)=\frac{\sum_iw_i^4}{(\sum_iw_i^2)^2}
 \ge\left(1+\frac{R(0.998)}{0.998^2}\right)^{-2}
 >0.8756>\frac78.
\]

For the physical purity at length $120$, return briefly to temperature $b_0$.
Its largest Boltzmann factor $z_0$ exceeds one by the trial bound, while the
sum of all factors is at most $q$.  After raising the factors to the power
$5/2$, the ratio of the remaining mass to $z_0^{5/2}$ is at most
$(q-1)^{5/2}<R_q$.  Therefore
\[
 T(120,b)\ge(1+R_q)^{-2}>0.8760>\frac78.
\]
This last step concerns physical Boltzmann factors, separately from the
spatial weights used to bound $S$.
\end{proof}

\subsection{The initialization giving the stronger gap}

\begin{proposition}\label{prop:seed2304}
At inverse temperature $b=784$,
\[
 1-S(2304,b)\le\frac{84513}{10^7}<\frac1{100},\qquad
 1-T(4608,b)\le\frac{27493}{5\cdot10^6}<\frac1{100}.
\]
\end{proposition}

\begin{proof}
We isolate one large eigenvalue in $J$ and bound the remaining moments.
This gives a pair of initial purity defects at length $36$; six rounded
updates then reach length $2304$.

Begin at $b_0=49/4$.  The $b_0=49/4$ table gives centers for $J_n,N_n$
through~\eqref{eq:sector-inversion}, with error at most
$\epsilon=30\cdot10^{-6}$.  Use the two filters
\[
 F_J(x)=x^4(7+17x-280x^2-185x^3+1000x^4)^2,
 \qquad F_N(x)=x^4(12-394x^2+1000x^4)^2.
\]
Their degree is $12$, so the table supplies every required moment.  The
summed upper bounds in Lemma~\ref{lem:polynomial-filter} are exactly
\[
 B_J=318604.54848175,\qquad B_N=48169.880699.
\]
Set $v_J=100139/100000$ and $v_N=872/1000$.  Direct rational comparisons give
\[
 F_J(v_J)-B_J>180,\quad F_J(-v_J)-B_J>496888,\quad
 F_N(\pm v_N)-B_N>654.
\]
To check the whole exterior rays without further numerical tests, write
$Q(x)=7+17x-280x^2-185x^3+1000x^4$.  For $y\ge0$,
\[
\begin{aligned}
 Q(1+y)&=559+2902y+5165y^2+3815y^3+1000y^4,\\
 Q(-1-y)&=895+3978y+6275y^2+4185y^3+1000y^4.
\end{aligned}
\]
Both are positive and increasing.  The polynomial
$12-394r^2+1000r^4$ is likewise positive and increasing for $r\ge v_N$:
its value at $v_N$ is positive and its derivative is
$4r(1000r^2-197)>0$.  Thus the filters increase along the requisite rays,
and every eigenvalue has modulus less than $v_J$ in $J$ and less than
$v_N$ in $N$.

The twelfth-moment upper bounds are
\[
 m_J=1.0657205,\qquad m_N=0.2783155.
\]
Put $\ell=999/1000$ and, for $k=36,72$, define
\begin{equation}\label{eq:seed72-tail}
 B(k)=3\mathcal C(m_N,v_N^{12},k/12),\qquad
 R_A(k)=B(k)+(m_J-\ell^{12})^{k/12}.
\end{equation}
Here $m_J>\ell^{12}$.  If all eigenvalues in $J$ had modulus at most $\ell$,
Lemma~\ref{lem:cappedmass} would imply
\[
 Z_{72}(b_0)\le\mathcal C(m_J,\ell^{12},6)+B(72)<1,
\]
where the strict margin to one exceeds $0.0693$.  The trial inequality
contradicts this.  Choose an eigenvalue $\lambda_0$ in $J$ with
$|\lambda_0|>\ell$.  Removing it leaves a twelfth-moment mass at most
$m_J-\ell^{12}$ in $J$.  The uncapped bound of Lemma~\ref{lem:cappedmass}
and the three $N$ sectors show that $R_A(k)$ bounds the remaining full
$k$th moment.

It follows that
\begin{equation}\label{eq:seed72-pq}
\begin{aligned}
 1-S(36,b_0)&\le p_0:=1-\left(1+\frac{R_A(36)}{\ell^{36}}\right)^{-2},\\
 1-T(72,b_0)&\le q_0:=1-\left(v_J^{72}+R_A(72)\right)^{-2}.
\end{aligned}
\end{equation}
For the first estimate retain $\lambda_0^{72}$ in the numerator; for the
second use $Z_{72}(2b_0)>1$ and
$Z_{72}(b_0)\le v_J^{72}+R_A(72)$.  Numerically $p_0<0.047916$ and
$q_0<0.181521$, but the following updates use their exact rational
definitions in~\eqref{eq:seed72-tail}--\eqref{eq:seed72-pq}.

Apply Lemma~\ref{lem:rounded-update} six times, rounding
each new defect upward to a multiple of $10^{-7}$ and using the rounded
spatial defect in the physical update.  For completeness, with
$f(u)=u^2/[2(1-u)^2]$ and $\lceil x\rceil_7=10^{-7}\lceil10^7x\rceil$,
the scalar updates are
\[
 p_{j+1}=\left\lceil f\bigl(1-(1-p_j)^2(1-q_j)\bigr)\right\rceil_7,
 \qquad
 q_{j+1}=\left\lceil f\bigl(1-(1-q_j)^2(1-p_{j+1})\bigr)\right\rceil_7.
\]
Exact rational evaluation gives
\[
\begin{array}{c|cc}
 j&p_j&q_j\\ \hline
 1&151249/2500000&433453/2500000\\
 2&343303/5000000&1632381/10000000\\
 3&44601/625000&361773/2500000\\
 4&632939/10000000&1055161/10000000\\
 5&375793/10000000&445941/10000000\\
 6&84513/10000000&27493/5000000
\end{array}
\]
Every update is in $[0,1)$, where $f$ is increasing, so the upward rounding
preserves valid bounds.  Both the length and inverse temperature have
multiplied by $2^6$: $36\cdot64=2304$ and $(49/4)\cdot64=784$.
The last row proves the proposition.
\end{proof}

\subsection{A periodic interface for boundary estimates}

The boundary-field companion~\cite{BoundaryCompanion} uses the intermediate spatial information
in the preceding proof, rather than only the final periodic gap.
We state that information separately so that its hypotheses and finite
inputs can be invoked together.

\begin{corollary}[Periodic inputs for boundary estimates]
\label{cor:boundary-inputs}
Let $X_b$ be the spatial transfer operator of
Proposition~\ref{prop:transfer}, with the shift $a=700741/500000$,
and put
\[
 \begin{gathered}
 b_0=49/4,\qquad b_1=49/2,\\
 \ell=999/1000,\qquad U=100139/100000,\qquad V=872/1000,\\
 m_J=1.0657205,\qquad m_N=.2783155.
 \end{gathered}
\]
At $b=b_0$, the full eigenvalue list, counted with multiplicity,
is the $D_2$-invariant list $J$ together with three identical lists $N$.
Every eigenvalue in $J$ has modulus less than $U$, and the sum of the
twelfth powers of all eigenvalues in $J$ is at most $m_J$.  Every eigenvalue in each $N$ list
has modulus less than $V$, and that list's twelfth powers sum to at
most $m_N$.  There is an eigenvalue in $J$ whose modulus exceeds $\ell$.
Moreover,
\[
 1-S(72,b_1)\le\frac{151249}{2500000}.
\]
For $n=72$ and $n=120$, one has $E_0(n)+na<0$; hence a physical ground
Boltzmann factor in $Z_n(b)$ exceeds one for every $b>0$.
\end{corollary}

\begin{proof}
The decomposition is Lemma~\ref{lem:sectors}, with $J=I+2O$.
The filters, twelfth-moment bounds and dominant eigenvalue are proved
at the start of the proof of Proposition~\ref{prop:seed2304}.
The first row of its rounded update table gives the displayed bound
at $(72,49/2)$.  The variational assertion is
Lemma~\ref{lem:trial-inputs}.
\end{proof}
